\section*{MỞ ĐẦU}
\phantomsection\addcontentsline{toc}{section}{\numberline {}MỞ ĐẦU}
Sự gia tăng nhanh chóng của Internet of Things đã tạo ra một nhu cầu lớn cho việc kết nối hàng tỷ thiết bị thông qua mạng không dây. Điều này đặt ra một loạt các thách thức về hiệu suất mạng và quản lý tài nguyên. Các mạng IoT thường sử dụng nguồn năng lượng có hạn, và việc sử dụng nguồn năng lượng tái tạo như năng lượng thu được từ môi trường (energy harvesting) trở nên quan trọng để giảm thiểu sự phụ thuộc vào nguồn năng lượng dự phòng. Hiệu suất mạng IoT cần phải tối ưu hóa để đảm bảo truyền thông chính xác và hiệu quả từ các thiết bị IoT đến trung tâm dữ liệu hoặc hệ thống giám sát.

Mục tiêu chính của đề tài là nghiên cứu và phát triển các phương pháp tối ưu hóa thông tin được truyền của các thiết bị IoT trong khi bị giới hạn bởi năng lượng, băng thông, và khả năng tính toán của chúng. Tạo ra một hệ thống thông minh có khả năng tự động điều chỉnh tốc độ cảm biến, chọn lựa phương pháp truyền dữ liệu và phân phối tài nguyên năng lượng dựa trên nguồn năng lượng hiện có và yêu cầu truyền thông của các thiết bị IoT. Xây dựng một hệ thống có khả năng tự động cân nhắc giữa hiệu suất mạng và tiết kiệm năng lượng.

Đối tượng nghiên cứu của đề tài là các mạng IoT sử dụng nguồn năng lượng tái tạo, đặc biệt là năng lượng thu được từ energy harvesting. Phạm vi nghiên cứu bao gồm việc phân tích, thiết kế, và triển khai các thuật toán để tối ưu hóa truyền thông và quản lý tài nguyên trong mạng IoT.

Đề tài đóng góp vào lĩnh vực nghiên cứu về mạng IoT, energy harvesting và tối ưu hóa tài nguyên. Nó cung cấp các phương pháp mới để nâng cao hiệu suất mạng IoT và giảm thiểu sự lãng phí nguồn năng lượng. Các kết quả của đề tài có thể được áp dụng trong các ứng dụng thực tế, như mạng cảm biến không dây cho quản lý môi trường, theo dõi y tế, quản lý năng lượng, và nhiều ứng dụng IoT khác, giúp tăng cường hiệu quả và tiết kiệm năng lượng.


\textbf{Từ khóa: }IoT, energy harvesting, Semcom, deep Q-network,…
\cleardoublepage